\chapter{Алгоритм поиска индуцированной активности с использованием нелинейного вложения (TriCo / eSPoC)}
\label{ch:trico_methodology}

%\textit{Теоретическое обоснование и вычислительный пайплайн, описываемые в данной главе, %были опубликованы соискателем в работах \cite{vakbib2, vakbib3}.}

\section{Концепция гибридного подхода}
\label{sec:trico_concept}

Как было показано в главе~\ref{ch:literature_review}, существует фундаментальный разрыв между высокой чувствительностью нелинейных графовых вложений (таких как UMAP) к сложной локальной динамике состояний мозга и строгой физической интерпретируемостью линейных методов пространственной фильтрации. 

Для преодоления этого ограничения в данной работе предлагается алгоритм Extended Source Power Comodulation (eSPoC), программно реализованный в виде пайплайна TriCo. Ключевая идея метода заключается в извлечении нелинейной низкоразмерной структуры непосредственно из ЭЭГ-данных и последующем использовании координат этого вложения в качестве многомерных регрессионных переменных (таргетов) для линейного пространственного фильтра.

\section{Вычислительный пайплайн алгоритма TriCo}
\label{sec:trico_algorithm}

Реализация метода состоит из трёх последовательных этапов: проекции в касательное пространство (Tangent space), нелинейного снижения размерности (UMAP) и линейной пространственной фильтрации (mSPoC).

\subsection{Проекция в касательное пространство (Tangent Space Mapping)}
На первом этапе непрерывный ЭЭГ-сигнал разбивается на $N_e$ временных эпох (окон), для каждой из которых вычисляется пространственная ковариационная матрица $\mathbf{C}_i \in \mathcal{S}_{++}^{N_c}$, где $i \in \{1, \dots, N_e\}$.

Поскольку матрицы ковариации лежат на Римановом многообразии, прямое применение к ним алгоритмов нелинейного вложения может быть субоптимальным. Для линеаризации пространства используется проецирование матриц в касательное пространство к многообразию в точке глобального риманова среднего $\mathbf{C}_{ref}$ \cite{yger2017riemann}:
\begin{equation}
	\mathbf{S}_i = \text{uvec}\left( \log \left( \mathbf{C}_{ref}^{-1/2} \mathbf{C}_i \mathbf{C}_{ref}^{-1/2} \right) \right),
\end{equation}
где оператор $\text{uvec}(\cdot)$ выполняет векторизацию верхней треугольной части симметричной матрицы, извлекая вектор признаков $\mathbf{s}_i \in \mathbb{R}^{N_c(N_c+1)/2}$. Матрица $\mathbf{S}_{feat} = [\mathbf{s}_1, \dots, \mathbf{s}_{N_e}]$ содержит евклидовы признаки для всех эпох.

\subsection{Нелинейное вложение (UMAP)}
Для извлечения глобальной поведенческой или когнитивной динамики (например, перехода между задачами или изменения состояния покоя) к полученной матрице признаков $\mathbf{S}_{feat}$ применяется алгоритм UMAP \cite{mcinnes2018umap}. 

Результатом работы алгоритма является низкоразмерное представление $\mathbf{Z} \in \mathbb{R}^{d \times N_e}$ (как правило, $d \in [2, 5]$). Векторы-столбцы $\mathbf{z}_i$ описывают эволюцию функционального состояния мозга на нелинейном манифолде с течением времени. 

\subsection{Оптимизация пространственного фильтра (mSPoC)}
\label{sec:trico_mspoc}
На финальном этапе координаты вложения $\mathbf{Z}$ принимаются в качестве многомерной внешней переменной для алгоритма mSPoC \cite{dahne2013mspoc}. Обозначим $\mathbf{w}$ как искомый пространственный фильтр для сенсорных данных $\mathbf{X}$, а $\mathbf{v}$ --- как весовой вектор для линейной комбинации осей UMAP-вложения.

Задача состоит в максимизации ковариации между мощностью отфильтрованного ЭЭГ сигнала и спроецированным значением на манифолде:
\begin{equation}
	(\mathbf{w}_{opt}, \mathbf{v}_{opt}) = \arg\max_{\mathbf{w}, \mathbf{v}} \frac{\mathbf{w}^T \mathbf{C}_{\mathbf{z}(\mathbf{v})} \mathbf{w}}{\mathbf{w}^T \mathbf{C} \mathbf{w}},
\end{equation}
где $\mathbf{C}$ --- усредненная матрица ковариации по всем эпохам, а взвешенная матрица $\mathbf{C}_{\mathbf{z}(\mathbf{v})}$ вычисляется как:
\begin{equation}
	\mathbf{C}_{\mathbf{z}(\mathbf{v})} = \frac{1}{N_e} \sum_{i=1}^{N_e} \left( \mathbf{v}^T \mathbf{z}_i - \mathbf{v}^T \bar{\mathbf{z}} \right) \mathbf{C}_i.
\end{equation}

Данная оптимизационная задача сводится к обобщенной задаче на собственные значения. В результате алгоритм глобально (через CCA-подобную логику) находит такое подпространство пространственных фильтров $\mathbf{W}$, которое наилучшим образом объясняет сложную траекторию состояний мозга, извлеченную UMAP.

\section{Дизайн симуляционных экспериментов}
\label{sec:trico_simulations}

Для проверки эффективности алгоритма TriCo был разработан ряд симуляций (toy examples). Модель генерации сигнала включала целевые нейрональные источники, амплитуда которых плавно комодулировала с заданными многомерными распределениями (имитируя реальную динамику на манифолде), а также сильный пространственный шум и фоновую активность (volume conduction). 

В экспериментах метод TriCo (eSPoC) напрямую сравнивался с базовыми жадными алгоритмами ($\text{SPoC}_\lambda$, mSPoC). 

% Далее в тексте будут описаны сами результаты симуляций с вставкой соответствующих графиков.